\documentclass[11pt,a4paper]{article}
\usepackage[T1]{fontenc}
\usepackage[utf8]{inputenc}
\usepackage{lmodern}
\usepackage[margin=2.6cm]{geometry}
\usepackage{microtype}
\usepackage{amsmath,amssymb,amsthm}
\usepackage{booktabs}
\usepackage{enumitem}
\usepackage[colorlinks,linkcolor=blue,citecolor=blue,urlcolor=blue]{hyperref}

\newcommand{\Qd}{\mathcal{Q}}
\DeclareMathOperator{\Tp}{Tp}
\theoremstyle{plain}
\newtheorem{observation}{Observation}

\title{\textbf{Toward $\Qd_8$: a feasibility study}}
\author{Chern++}
\date{September 2026}

\begin{document}
\maketitle

\begin{abstract}
The Morin Thom polynomials are computed from a single polynomial $\Qd_d$, the equivariant
multidegree of a Borel orbit closure. $\Qd_d$ is known for $d \le 7$, and $d = 8$ is out of reach
of Gr\"obner-basis methods. A formula for $\Qd_8$ would test the conjectures of
\texttt{report/ballot.pdf} at $d = 8$. We test the natural resolution candidate, the rank-one
nonassociative tower, which has an explicit product formula for its point class. It agrees
with $\Qd_4$ but differs from $\Qd_d$ in the residue-relevant sense for $d = 5, 6, 7$. Its error
is a product of linear forms at $d = 5$ and then becomes irreducible: a sextic times a linear form
at $d = 6$, and an irreducible polynomial of degree $13$ at $d = 7$. So this route does not
reach $\Qd_8$, and no other resolution is presently in hand.
\end{abstract}

\section{Setting}

Let $N_d \subset \mathrm{Hom}(\mathbb{C}^d, \mathrm{Sym}^2\mathbb{C}^d)$ be spanned by the coordinates
$q^{mr}_l$ with $1 \le m \le r$ and $m + r \le l \le d$. The coordinate $q^{mr}_l$ has torus weight
$z_m + z_r - z_l$. The reference point $\varepsilon$ has $q^{mr}_l = 1$ exactly when $m + r = l$: it
is the multiplication table of $t\mathbb{C}[t]/(t^{d+1})$ in the basis $t, \dots, t^d$. The
upper-triangular Borel subgroup $B_d \subset GL_d$ acts on $N_d$. Then
\[
  \Qd_d = [\overline{B_d \cdot \varepsilon}] \in \mathbb{Z}[z_1, \dots, z_d],
\]
the equivariant multidegree of the orbit closure. It is homogeneous of degree
$\dim N_d - \binom d2 = 1, 3, 7, 13, 22$ for $d = 4, \dots, 8$. By the B\'erczi--Szenes residue formula
it determines $\Tp_{A_d}$ at every relative dimension. It matters only modulo the
\emph{residue-null} polynomials, those that do not change any Chern coefficient. In this project
$\Qd_4, \dots, \Qd_7$ are computed by Gr\"obner bases (SageMath/Singular) and stored as artifacts.
The cost of the analogous $d = 8$ computation was estimated at $\sim 10^{11}$ times that of $d = 7$.

\section{The candidate: the nonassociative tower}

A resolution of the orbit closure would give $\Qd_d$ by equivariant localisation. The natural
candidate is the tower of \emph{rank-one nonassociative} structures, built level by level. Its
point class is an explicit product of weights:
\[
  P^\star_d = \prod_{k=4}^{d} \ \prod_{\substack{2 \le i \le j \\ i + j \le k}} (z_1 + z_{i-1} + z_j - z_k).
\]
The number of factors is $1, 3, 7, 13, 22$ for $d = 4, \dots, 8$, matching $\deg \Qd_d$. This class and
the observation that it recovers $\Qd_4$ are due to an unrefereed external
report~\cite{FS}. Everything below is recomputed here from the formula above and the stored
$\Qd_d$.

\begin{observation}[verified]\label{obs}
\begin{enumerate}[nosep]
\item $P^\star_4 = \Qd_4$.
\item $\Qd_d - P^\star_d$ is not residue-null for $d = 5, 6, 7$. Replacing $\Qd_d$ by $P^\star_d$
  changes $21$ of the $32$ testable packet sums at $d = 5$, $11$ of $13$ at $d = 6$, and $1$ of $3$ at
  $d = 7$, on the truncations tested.
\item The differences factor over $\mathbb{Q}$ as follows:
  \begin{itemize}[nosep]
  \item at $d = 5$, $\Qd_5 - P^\star_5 = (2z_1 - z_2)(z_1 + z_3 - z_4)(2z_1 + z_2 - z_5)$;
  \item at $d = 6$, a linear form $(2z_1 - z_2)$ times an irreducible sextic ($395$ terms in all);
  \item at $d = 7$, an irreducible polynomial of degree $13$ ($14{,}552$ terms).
  \end{itemize}
\end{enumerate}
\end{observation}

\paragraph{Reading.} $P^\star_d$ is the class of a variety that contains the orbit closure and,
from $d = 5$, has extra components. So the tower is not birational onto
$\overline{B_d\varepsilon}$. At $d = 5$ the defect is a single product of linear forms. That is
consistent with one extra linear component, and it suggests writing $\Qd_d$ as $P^\star_d$ minus
component classes. From $d = 6$ the defect has an irreducible factor of high degree, which rules
out a correction by a few linear components. At $d = 7$ nothing factors at all. The defect grows
roughly like $\Qd_d$ itself, so the tower does not reduce the problem.

\section{What a route to $\Qd_8$ would need}

\begin{itemize}[nosep]
\item \textbf{A genuinely birational smooth model of $\overline{B_d\varepsilon}$}, compatible
  with the torus, with isolated fixed points. Localisation would then give $\Qd_d$ as a sum over
  fixed points. The tower above fails this. The saturated associativity scheme has the right
  image but is singular.
\item \textbf{Or a characterisation of $\Qd_d$ modulo the residue-null kernel only.} Only the
  packet class of $\Qd_d$ enters the Thom polynomial. Several constraints already pin that class
  down partially:
  \begin{itemize}[nosep]
  \item the degree;
  \item the normalisation;
  \item stabilisation, $[z_d^{e(d)}]\Qd_d = \pm \Qd_{d-1}$ with $e(d) = \lfloor d^2/4 \rfloor - d + 1$
    (verified for $d \le 7$);
  \item the plane-sector theorem, which fixes every packet with $\max M \le 1$ (given its two
    stated geometric inputs);
  \item Rim\'anyi's published $\Tp_{A_8}$ at $\ell = 0$.
  \end{itemize}
  Whether these determine the packet class of $\Qd_8$ is open. Counting constraints against the
  dimension of the space of admissible numerators is the natural next computation.
\end{itemize}

\paragraph{Decision.} Following the gate in \texttt{report/research\_plan.md}, the $\Qd_8$ computation
is not attempted. The d = 8 question that matters most, $\kappa_8 < 1$ in the dominance conjecture
and the ballot inequality at $d = 8$, stays open.

\paragraph{Reproducing.} The comparison uses \texttt{chernpp.optimisation.gauge.validate\_gauge} on
the kernel $\Qd_d - P^\star_d$, and Singular's \texttt{factorize} on the difference.

\begin{thebibliography}{9}
\bibitem{FS} Rim\'anyi positivity: findings summary (29 July 2026), unrefereed, \S8.2. Cited for the tower and the formula for $P^\star_d$ only; every claim above is recomputed.
\end{thebibliography}
\end{document}
